% conclusion.tex
\section{Conclusion}\label{concl}
This paper examines a novel approach to uncovering anti-competitive behavior in public procurement, applying it to a large dataset of auctions from São Paulo, Brazil. By focusing on auctions decided by very small margins, we are able to detect subtle yet systematic differences between firms that narrowly win and those that narrowly lose. Our regression discontinuity design reveals that narrowly winning firms are far more likely to have won previous contracts and tend to carry a heavier workload of ongoing projects than their narrowly losing counterparts. This evidence points to the presence of an incumbency advantage in the bidding process that is inconsistent with purely competitive behavior.

The weight of the evidence suggests that this advantage is driven by preferential treatment or corrupt favoritism rather than a classic collusive arrangement among firms. We find no indication of the kind of bid-rotation patterns one would expect if a cartel were in operation; instead, the data show a concentration of wins and contracts in the hands of specific firms. In the context of Brazilian public procurement, this could reflect political connections or bribery influencing contract awards, whereby certain firms become the "favorites" that win disproportionately often. Our contribution is to provide a data-driven method to flag such favoritism: the discontinuities in incumbency and backlog at the win/loss threshold serve as telltale signs of a distorted playing field.

These findings have important policy implications. Ensuring fairness in public procurement is crucial for efficiency and public trust. The methodology and results we present can help procurement authorities and regulators identify potential problems early. For instance, procurement oversight bodies could incorporate threshold-based analysis into their monitoring systems to highlight cases where a firm’s winning record is abnormally high given close competition. Armed with such tools, auditors or anti-corruption agencies could target investigations toward agencies or contracts that exhibit the red flags consistent with favoritism or collusion. Additionally, reforms that increase transparency—such as random audits, public release of detailed bidding data, and stricter enforcement of conflict-of-interest rules—may help deter the kind of behavior we document by raising the risk of detection.

While our analysis provides strong evidence of favoritism in São Paulo’s procurement auctions, it also opens avenues for further research. One limitation of our current approach is that it infers collusion or favoritism indirectly from outcomes; future work could combine this with direct evidence (such as investigation records or social network links between contractors and officials) to conclusively identify culpable parties. Moreover, the application of our method to other contexts (other regions or countries, and different procurement systems) would be valuable to see how widespread such issues are and whether similar patterns emerge under different institutional settings. 

In conclusion, competitive procurement is fundamental for obtaining value for public money, and our study shows how modern econometric tools can peer beneath the surface of auction data to reveal when that competition is only an illusion. By identifying concrete signals of anti-competitive behavior, we hope to contribute to more accountable and efficient public contracting processes in Brazil and beyond.
